% Generated draft: requires mathematical and source review.
\documentclass[11pt]{article}
\usepackage[margin=25mm]{geometry}
\usepackage{fontspec}
\setmainfont{Latin Modern Roman}
\usepackage{amsmath,amssymb,mathtools,unicode-math}
\setmathfont{Latin Modern Math}
\usepackage{enumitem,needspace,xurl,hyperref}
\setlength{\parindent}{0pt}
\setlength{\parskip}{5pt}
\setlength{\emergencystretch}{3em}
\newcommand{\N}{\mathbb{N}}
\newcommand{\Z}{\mathbb{Z}}
\newcommand{\Q}{\mathbb{Q}}
\newcommand{\R}{\mathbb{R}}
\newcommand{\C}{\mathbb{C}}
\newcommand{\F}{\mathbb{F}}
\newcommand{\A}{\mathbb{A}}
\newcommand{\PP}{\mathbb{P}}
\newcommand{\E}{\mathbb{E}}
\newcommand{\eps}{\varepsilon}
\newcommand{\dd}{\,\mathrm{d}}
\begin{document}
\section*{UnsolvedMath ID: 30000219}
\textbf{Source problem code:} catalogue entry
\subsection{Definitions and mathematical statement}
Let \(S=K[t_{ij}]\) be the polynomial ring of the source, grouped into row blocks indexed by \(1\leq i\leq m\), and give it the multigrading \(\deg(t_{ij})=e_i\in\Z^m\). Write \(L=L(a)\) for the intended matrix of linear forms with coefficients \(a_{ijk}\), and let \(I_2(L)\subset S\) be the ideal of its \(2\times2\) minors. The formula for \(L(a)\), matrix dimensions, and ranges of \(j,k\) are missing from the supplied record, so the allowed matrix family is not silently taken to be all matrices of linear forms.

For a monomial term order \(\prec\), the conjecture is that, for every allowed coefficient choice \(a\) and every \(\prec\), the ideal \(\operatorname{in}_{\prec}I_2(L(a))\) can be generated by monomials
\[
t_{i_1j_1}\cdots t_{i_sj_s},\qquad i_1<i_2<\cdots<i_s.
\]
Equivalently, every minimal monomial generator has multidegree in \(\{0,1\}^m\). This row-block squarefreeness is stronger than ordinary squarefreeness: distinct variables from the same row block cannot occur together in such a generator. The quantifier is every term order in the given variables, not existence of a favorable order or generic coordinate change.

The second question asks whether the same statement extends to every distributive lattice in the intended analogous construction. A distributive lattice has meet and join satisfying the distributive laws. The source does not give its associated polynomial ring, defining or deformed ideal, coefficient parameters, multigrading, or intended finiteness assumptions. Therefore no particular Hibi or straightening-law construction is substituted for the unspecified extension.

The supplied assessment checked 2026-08-21 reports the squarefree initial-ideal assertion proved for generic matrices, together with associated Cohen--Macaulay and Koszul consequences, but says that extension to arbitrary distributive lattices remains open in the checked literature. It does not reproduce the exact generic-matrix theorem or identify which coefficient specializations it covers. Those consequences are reported for the indicated established setting, not inferred solely from squarefreeness for all unspecified ideals.
\subsection{Short English statement}
Ask whether every term order and every allowed coefficient specialization yield initial generators using at most one variable from each row block, and whether the corresponding assertion extends to distributive lattices. The original matrix and lattice constructions are missing.
\subsection{Sources}
\begin{itemize}
\item Local UnsolvedMath registry, permanent ID 30000219, catalogue entry, \emph{Square-Free Initial Ideals of Linear Determinantal Ideals}.
\item \emph{Kommutative Algebra} (2006), Oberwolfach Reports, pp. 1073--1126, DOI \url{https://doi.org/10.4171/owr/2005/19}.
\item Supplied literature triage checked 2026-08-21 cites only the original report. No separate authoritative theorem reference or full definition of the matrix and lattice constructions is supplied.
\end{itemize}
\end{document}
